\documentclass[10pt]{article}
\usepackage[a4paper,margin=0.86in]{geometry}
\renewcommand{\rmdefault}{cmr}
\usepackage{amsmath,amssymb,amsthm,mathtools,amscd,enumitem,xcolor,hyperref}
\hypersetup{colorlinks=true,linkcolor=blue!55!black,citecolor=blue!55!black,urlcolor=blue!55!black}
\setlength{\emergencystretch}{4em}
\setlength{\parindent}{0pt}
\setlength{\parskip}{0.52em}
\setlist{leftmargin=1.55em,itemsep=0.2em,topsep=0.25em}
\sloppy

\newtheorem{theorem}{Theorem}[section]
\newtheorem{proposition}[theorem]{Proposition}
\newtheorem{lemma}[theorem]{Lemma}
\newtheorem{corollary}[theorem]{Corollary}
\theoremstyle{definition}
\newtheorem{definition}[theorem]{Definition}
\newtheorem{convention}[theorem]{Convention}
\newtheorem{remark}[theorem]{Remark}

\newcommand{\Zl}{\mathbf Z_\ell}
\newcommand{\Ql}{\mathbf Q_\ell}
\newcommand{\Qlb}{\overline{\mathbf Q}_\ell}
\newcommand{\OO}{\mathcal O}
\newcommand{\Dbc}{D^b_c}
\newcommand{\Dctf}{D^b_{\mathrm{ctf}}}
\newcommand{\RHom}{R\mathcal Hom}
\newcommand{\red}{\operatorname{red}}
\newcommand{\DD}{\mathbb D}
\newcommand{\Tr}{\operatorname{Tr}}

\begin{document}
\begin{center}
{\large\bfseries Verdier Duality and Relative Duality\par}
\medskip
{\normalsize Dualizing complexes, biduality, and exchange morphisms\par}
\end{center}

\begin{abstract}
This note constructs Verdier duality from the derived internal Hom and the
exceptional inverse image of a base dualizing object.  Constructibility,
biduality, open and proper exchanges, smooth purity, relative duality,
pairings, and coefficient-change rectangles are stated as a sequence of
separate results.  All source and target categories are explicit.
\end{abstract}

\tableofcontents

\section{Coefficient stages and the six operations used below}

Fix a prime \(\ell\).  Schemes are noetherian of finite Krull dimension,
\(\ell\) is invertible on them, and morphisms are quasi-compact and
quasi-separated.  Sheaves are taken on the small pro-etale site.

\begin{definition}[Finite coefficient ambient category]
Let \(\Lambda\) be a finite local
\(\mathbf Z/\ell^n\mathbf Z\)-algebra and let \(\Lambda_X\) be the constant
sheaf of rings on \(X_{\mathrm{pro\acute et}}\).  Write

\[
 \mathcal A_X(\Lambda)
 =\mathsf{Mod}(X_{\mathrm{pro\acute et}},\Lambda_X),
 \qquad
 D(X,\Lambda)=D\bigl(\mathcal A_X(\Lambda)\bigr),
\]

where the latter is the Verdier localization of cochain complexes at
quasi-isomorphisms.  The full triangulated subcategory
\(\Dbc(X,\Lambda)\) consists of objects having finitely many nonzero
cohomology sheaves, each locally constant with finite stalks on every member
of some finite constructible stratification of \(X\).  The subcategory
\(\Dctf(X,\Lambda)\) consists of those \(K\in\Dbc(X,\Lambda)\) for which
there are integers \(r\leq s\) such that

\[
 K\otimes^L_{\Lambda_X}M\in D^{[r,s]}(X,\Lambda)
\]

for every sheaf \(M\) of \(\Lambda_X\)-modules placed in degree zero.
\end{definition}

\begin{definition}[Normalized adic ambient category]
Let \(E/\Ql\) be finite, let \(\OO_E\) be its ring of integers, choose a
uniformizer \(\varpi\), and put \(\Lambda_n=\OO_E/\varpi^n\).  Let
\(\mathcal A_X(\Lambda_\bullet)\) be the abelian category of inverse systems
\(K_\bullet=(K_n,t_{n+1,n})\), where \(K_n\) is a sheaf of
\(\Lambda_{n,X}\)-modules and

\[
 t_{n+1,n}:K_{n+1}\longrightarrow K_n
\]

is linear through \(\Lambda_{n+1}\to\Lambda_n\).  Put

\[
 D_{\mathrm{sys}}(X,\OO_E)
 =D\bigl(\mathcal A_X(\Lambda_\bullet)\bigr).
\]

Evaluation at level \(n\) is an exact triangulated functor
\(\operatorname{ev}_{n,X}:D_{\mathrm{sys}}(X,\OO_E)\to D(X,\Lambda_n)\).
An object \(K\) is \emph{normalized} if, for every \(m\leq n\), the
transition-induced morphism

\[
 \nu^K_{n,m}:
 \operatorname{ev}_{n,X}K\otimes^L_{\Lambda_{n,X}}\Lambda_{m,X}
 \xrightarrow{\sim}\operatorname{ev}_{m,X}K
\]

is an isomorphism, with \(\nu^K_{n,n}=\operatorname{id}\) and the usual
reassociation identity for \(r\leq m\leq n\).  Define

\[
 D(X,\OO_E)=D_{\mathrm{nor}}(X,\OO_E)
 \subset D_{\mathrm{sys}}(X,\OO_E)
\]

to be the full triangulated subcategory of normalized objects.  Define
\(\Dbc(X,\OO_E)\) by requiring the objects
\(K_n=\operatorname{ev}_{n,X}K\) to have one common cohomological interval
and one compatible finite constructible stratification.  Define
\(\Dctf(X,\OO_E)\) by additionally requiring one common Tor-amplitude
interval.
\end{definition}

\begin{definition}[Finite coefficient-field ambient category]
For the chosen lattice \(\OO_E\), let
\(D(X,\OO_E)[1/\varpi]\) have the same objects as
\(D(X,\OO_E)\) and morphism groups

\[
 \operatorname{Hom}_{D(X,\OO_E)[1/\varpi]}(K,L)
 =\operatorname{Hom}_{D(X,\OO_E)}(K,L)\otimes_{\OO_E}E.
\]

Define

\[
 D(X,E)=\operatorname{Kar}\bigl(D(X,\OO_E)[1/\varpi]\bigr).
\]

Thus an object is a pair \((K,e)\), with \(e^2=e\) in the localized
endomorphism ring, and a morphism \((K,e)\to(L,e')\) is a localized
morphism \(u:K\to L\) satisfying \(u=e'ue\).  The categories
\(\Dbc(X,E)\) and \(\Dctf(X,E)\) are the thick idempotent-complete
subcategories represented by \(\Dbc(X,\OO_E)\) and
\(\Dctf(X,\OO_E)\), respectively.  Scalar extension along \(E\to F\)
is induced by extension of lattices and is independent of the chosen
lattices up to its transitive canonical isomorphism.
\end{definition}

\begin{definition}[Algebraically closed coefficient ambient category]
Let \(\mathcal C_X\) be the ordinary \(\Ql\)-linear category whose objects
are pairs \((E,K)\), with \(E/\Ql\) finite and \(K\in D(X,E)\), and whose
morphism groups are

\[
 \operatorname{Hom}_{\mathcal C_X}((E,K),(E',L))
 =\varinjlim_{F\supset E,E'}
 \operatorname{Hom}_{D(X,F)}(K\otimes_EF,L\otimes_{E'}F).
\]

The limit runs over finite extensions \(F/\Ql\) containing images of both
fields; composition is composition after passage to one common finite
extension.  Define

\[
 D(X,\Qlb)=\operatorname{Kar}(\mathcal C_X).
\]

An object is therefore a triple \((E,K,e)\), up to finite scalar extension,
where \(e\) is an idempotent endomorphism.  The full thick subcategories
\(\Dbc(X,\Qlb)\) and \(\Dctf(X,\Qlb)\) consist of triples represented by
the corresponding subcategories over a finite \(E\).
\end{definition}

\begin{convention}[Admissible coefficient symbol]
Only after the preceding four definitions, let \(A\) denote one of
\(\Lambda\), \(\OO_E\), \(E\), or \(\Qlb\).  The notation \(D(X,A)\),
\(\Dbc(X,A)\), and \(\Dctf(X,A)\) always means the category defined at
that coefficient stage.  The symbol \(A_X\) means \(\Lambda_X\) at a finite
stage, the normalized constant system \((\Lambda_{n,X})_n\) at the integral
stage, and its scalar extension or finite coefficient-field representative at the
\(E\)- or \(\Qlb\)-stage.  A formula over \(E\) or \(\Qlb\) is asserted only
after the corresponding finite-stage functor and scalar-extension comparison
have been stated.
\end{convention}

\begin{definition}[Ordinary inverse and direct image]
For finite \(A=\Lambda\), the morphism of ringed pro-etale topoi associated
with \(f:X\to Y\) gives

\[
 f^*M=f^{-1}M\otimes_{f^{-1}\Lambda_Y}\Lambda_X,
 \qquad
 f_*N(U)=N(U\times_YX).
\]

Choose a functorial K-flat replacement \(P(K)\to K\) and a functorial
K-injective replacement \(L\to I(L)\).  Define

\[
 Lf^*K=f^*P(K):D(Y,\Lambda)\longrightarrow D(X,\Lambda),
\]

\[
 Rf_*L=f_*I(L):D(X,\Lambda)\longrightarrow D(Y,\Lambda).
\]

The underived tensor--Hom adjunction induces
\(Lf^*\dashv Rf_*\).  At the \(\OO_E\)-stage the same formulas are first
formed in
\(D_{\mathrm{sys}}(Y,\OO_E)\rightleftarrows
D_{\mathrm{sys}}(X,\OO_E)\).  Their finite-reduction comparisons give the
displayed functors on normalized objects whenever normalization is preserved.
The \(E\)- and \(\Qlb\)-functors are induced by rationalization, finite
scalar extension, and passage to idempotent summands.
\end{definition}

\begin{theorem}[Coefficient-stage landing of inverse and direct image]
The functor \(Lf^*\) preserves normalized systems and finite Tor-amplitude.
If the bounded constructible direct-image theorem and its finite-reduction
comparison hold for \(f\), then the system-level \(Rf_*\) preserves
normalization and restricts to

\[
 Rf_*:\Dbc(X,\OO_E)\longrightarrow\Dbc(Y,\OO_E).
\]

Under the corresponding uniform Tor-amplitude theorem it further restricts
to \(\Dctf\).  These restrictions commute with rationalization and finite
scalar extension and therefore define the same functors over \(E\) and
\(\Qlb\), with the source and target categories of the admissible coefficient
convention.
\end{theorem}

\begin{definition}[Compact-support direct image]
If \(f\) is separated and of finite type, choose a compactification
\(X\xrightarrow j\overline X\xrightarrow pY\) and define

\[
 Rf_!=Rp_*j_!:D(X,A)\longrightarrow D(Y,A).
\]

Common proper refinements give canonical comparison isomorphisms satisfying
the identity and cocycle equations.
\end{definition}

\begin{definition}[Full site-local cohomological-dimension datum]
Choose a set \(\mathcal B_Y\) of affine, quasi-compact, weakly contractible
objects covering the selected small pro-etale site of \(Y\).  For a proper
map \(p:Z\to Y\), the assertion
\(\operatorname{FullCD}(p,A;c)\) consists of one integer \(c\) such that

\[
 R^qp_*F=0\qquad(q>c)
\]

for every sheaf \(F\) of \(A\)-modules on \(Z_{\mathrm{pro\acute et}}\), and
such that the same integer works after every selected site-local base
change.  Explicitly, for \(V\in\mathcal B_Y\), put
\(p_V:Z\times_YV\to V\); then

\[
 R^q(p_V)_*G=0\qquad(q>c)
\]

for every sheaf \(G\) of \(A\)-modules on
\((Z\times_YV)_{\mathrm{pro\acute et}}\).  A bound only for constructible
sheaves is not a \(\operatorname{FullCD}\) datum.
\end{definition}

\begin{theorem}[Ambient exceptional adjoint at finite coefficients]
At a finite coefficient stage, assume
\(\operatorname{FullCD}(p,A;c)\) for the proper part \(p\) of the chosen
compactification, and assume the weakly contractible compact-generator,
left-completeness, and arbitrary-coproduct theorems.  Then \(Rf_!\) preserves
arbitrary coproducts and has an exact right adjoint

\[
 Rf^!:D(Y,A)\longrightarrow D(X,A).
\]

with adjunction

\[
 Rf_!\dashv Rf^!.
\]
\end{theorem}

\begin{definition}[Finite reduction]
For \(m\leq n\), derived extension of scalars defines

\[
 \red_{n,m,X}:D(X,\Lambda_n)\longrightarrow D(X,\Lambda_m),
 \qquad
 \red_{n,m,X}(K)=K\otimes^L_{\Lambda_{n,X}}\Lambda_{m,X}.
\]

It restricts to \(\Dctf(X,\Lambda_n)\to\Dctf(X,\Lambda_m)\).  The canonical
associativity maps
\(\red_{m,r,X}\red_{n,m,X}\cong\red_{n,r,X}\) satisfy the identity and
cocycle equations.
\end{definition}

\begin{theorem}[Constructible and finite-Tor landing of the six operations]
Assume the finite compact-support amplitude and constructible-devissage
theorems for \(f\).  Then \(Rf_!\) and the Brown representing object \(Rf^!L\)
are bounded constructible for bounded constructible input.  If the
stratified finite-Tor theorem is also available, the adjunction restricts to

\[
 Rf_!: \Dctf(X,A)\rightleftarrows\Dctf(Y,A):Rf^!.
\]

Independently, the constructible direct-image theorem gives
\(Rf_*:\Dbc(X,A)\to\Dbc(Y,A)\), and its finite-Tor form gives
\(Rf_*:\Dctf(X,A)\to\Dctf(Y,A)\).  If \(f\) is proper, the compactification
comparison is \(Rf_!\cong Rf_*\); if \(j\) is open, adjoint uniqueness gives
\(Rj^!\cong j^*\).
\end{theorem}

\begin{convention}[Exchange hypotheses]
On the displayed finite-Tor categories, the projection formula, proper base
change for \(Rf_*\), compact-support base change for \(Rf_!\), finite
reduction, and exceptional scalar extension are separate comparison
theorems.  Each is required to satisfy its stated identity, composition, and
pasting equations.  None is inferred merely from the existence of \(Rf^!\).
\end{convention}

\begin{convention}[Scope of the six operations below]
Every use of \(Rf^!\) below assumes ambient exceptional representability and
the relevant constructibility and finite-Tor preservation theorem.  Every
displayed projection, base-change, reduction, or scalar-extension
isomorphism assumes its separately named exchange theorem.
\end{convention}

\section{Underived and derived internal Hom}

\begin{definition}[Sheaf internal Hom]
At a finite stage \(A=\Lambda\), for sheaves \(M,N\) of
\(\Lambda_X\)-modules let
\(\mathcal Hom_X(M,N)\) be the sheafification of

\[
 U\longmapsto
 \operatorname{Hom}_{\Lambda_U}(M|_U,N|_U).
\]

It is contravariant and left exact in \(M\), covariant and left exact in
\(N\), and satisfies

\[
 \operatorname{Hom}_X(L\otimes_{\Lambda_X}M,N)
 \cong
 \operatorname{Hom}_X(L,\mathcal Hom_X(M,N)).
\]
\end{definition}

\begin{definition}[Internal Hom complex]
For cochain complexes \(P\) and \(I\), define

\[
 \mathcal Hom_X^n(P,I)=
 \prod_{r\in\mathbf Z}\mathcal Hom_X(P^r,I^{r+n}).
\]

For a homogeneous section \(\phi\) of degree \(n\), set

\[
 d(\phi)=d_I\circ\phi-(-1)^n\phi\circ d_P.
\]

Then \(d^2=0\).  Evaluation is the cochain map

\[
 P\otimes_{A_X}\mathcal Hom_X^\bullet(P,I)\longrightarrow I,
 \qquad p\otimes\phi\longmapsto(-1)^{|p||\phi|}\phi(p),
\]

with the sign fixed by the symmetric monoidal convention for cochain
complexes.
\end{definition}

\begin{definition}[Derived internal Hom]
At a finite stage \(A=\Lambda\), choose a functorial K-flat replacement
\(P(K)\to K\) and a functorial
K-injective replacement \(L\to I(L)\).  Define

\[
 \RHom_X(K,L)=\mathcal Hom_X^\bullet(P(K),I(L))
 \in D(X,A).
\]

This gives an exact functor in each triangulated variable,

\[
 \RHom_{X,\Lambda}(-,-):D(X,\Lambda)^{\mathrm{op}}
 \times D(X,\Lambda)\longrightarrow D(X,\Lambda).
\]
\end{definition}

\begin{theorem}[Resolution independence]
Any two K-flat source replacements and K-injective target replacements give
canonically isomorphic internal Hom objects.  The comparisons satisfy
identity, composition, and cocycle equations; hence the preceding definition
is a functor on the Verdier localization.
\end{theorem}

\begin{theorem}[Derived tensor--Hom adjunction]
For \(K,L,M\in D(X,\Lambda)\), there is a natural isomorphism

\[
 \operatorname{Hom}_X(K\otimes^L_{\Lambda_X}L,M)
 \xrightarrow{\sim}
 \operatorname{Hom}_X(K,\RHom_X(L,M)).
\]

Its counit is the derived evaluation map

\[
 \operatorname{ev}_{L,M}:
 L\otimes^L\RHom_X(L,M)\longrightarrow M.
\]

 Currying followed by uncurrying and uncurrying followed by currying are the
 identity.
\end{theorem}

\begin{theorem}[Internal Hom at the adic and field stages]
Let \(K\in\Dctf(X,\OO_E)\) and \(L\in\Dbc(X,\OO_E)\), with finite
levels \(K_n\) and \(L_n\).  The finite-level objects

\[
 \RHom_{X,\Lambda_n}(K_n,L_n)
\]

and the finite-reduction isomorphisms

\[
 \red_{n,m,X}\RHom_{X,\Lambda_n}(K_n,L_n)
 \xrightarrow{\sim}
 \RHom_{X,\Lambda_m}(\red_{n,m,X}K_n,\red_{n,m,X}L_n)
\]

form a normalized object of \(D(X,\OO_E)\).  This object is, by definition,
\(\RHom_{X,\OO_E}(K,L)\).  Its evaluation map and tensor--Hom adjunction are
the maps whose reductions are the finite-level maps.  Rationalization defines
\(\RHom_{X,E}\).  For triples over \(\Qlb\), both variables are represented
over one common finite coefficient field; internal Hom reverses the idempotent in the
first variable and retains the idempotent in the second.  The resulting
object is independent of the common finite coefficient field and is transitive under
scalar extension.
\end{theorem}

\begin{convention}[Coefficient-qualified internal Hom]
After the preceding theorem, \(\RHom_{X,A}\), or \(\RHom_X\) when the
coefficient is visible, denotes the finite construction or its adic,
rational, or \(\Qlb\) reconstruction.  On bounded constructible categories
it is used with a finite-Tor first variable.
\end{convention}

\begin{definition}[Perfect object]
An object \(P\in D(X,A)\) is perfect if it is locally represented by a
bounded complex of finite projective \(A_X\)-modules.  Its coefficient dual
is

\[
 P^\vee=\RHom_X(P,A_X).
\]
\end{definition}

\begin{proposition}[Dualizability of perfect objects]
For a perfect \(P\), the canonical morphism

\[
 P^\vee\otimes^L K\xrightarrow{\sim}\RHom_X(P,K)
\]

is an isomorphism for every \(K\in D(X,A)\).  The evaluation and coevaluation
maps make \(P^\vee\) a two-sided dual of \(P\).  Conversely, every dualizable
object is perfect.
\end{proposition}

\begin{theorem}[Constructible closure of internal Hom]
If \(K\in\Dctf(X,A)\) and \(L\in\Dbc(X,A)\), then

\[
 \RHom_X(K,L)\in\Dbc(X,A).
\]

If both \(K\) and \(L\) have finite Tor-amplitude, then so does
\(\RHom_X(K,L)\).  Thus derived internal Hom restricts to the indicated
bounded constructible categories.
\end{theorem}

\section{Dualizing objects and the Verdier functor}

\begin{definition}[Dualizing object at one coefficient stage]
For one admissible coefficient \(A\), an object
\(\omega_{S,A}\in\Dbc(S,A)\) is dualizing on the finite-Tor constructible
category if

\[
 \RHom_{S,A}(-,\omega_{S,A}):
 \Dctf(S,A)^{\mathrm{op}}\longrightarrow\Dctf(S,A)
\]

and, for every \(K\in\Dctf(S,A)\), the evaluation-adjoint map

\[
 K\longrightarrow
 \RHom_{S,A}\bigl(\RHom_{S,A}(K,\omega_{S,A}),\omega_{S,A}\bigr)
\]

is an isomorphism.
\end{definition}

\begin{definition}[Compatible base dualizing objects]
A compatible coefficient family on \(S\) consists first of dualizing objects
\(\omega_{S,n}\in\Dctf(S,\Lambda_n)\) and isomorphisms

\[
 \lambda_{n,m}:\red_{n,m,S}\omega_{S,n}
 \xrightarrow{\sim}\omega_{S,m}\qquad(m\leq n)
\]

with \(\lambda_{n,n}=\operatorname{id}\) and

\[
 \lambda_{n,r}=\lambda_{m,r}\circ
 \red_{m,r,S}(\lambda_{n,m})
\]

after reassociation.  These data define a normalized object
\(\omega_{S,\OO_E}\in\Dctf(S,\OO_E)\).  Its rationalization defines
\(\omega_{S,E}\).  A choice over \(\Qlb\) is a triple
\((E,\omega_{S,E},e)\) in the pair/Karoubi category, and every further
finite scalar extension carries it to the specified equivalent triple.
Fix one such compatible family.  When \(S=\operatorname{Spec}k\) with
\(k\) separably closed, the standard family is the constant family
\(\omega_{S,A}=A_S\).
\end{definition}

\begin{definition}[Relative and absolute notation for dualizing objects]
For a separated finite-type structural morphism \(a_X:X\to S\), define at
each coefficient stage

\[
 \omega_{X/S,A}=Ra_X^!\omega_{S,A}\in D(X,A).
\]

Composition of exceptional inverse images gives, for
\(f:X\to Y\) over \(S\), a canonical isomorphism

\[
 \kappa_{f,A}:\omega_{X/S,A}\xrightarrow{\sim}Rf^!\omega_{Y/S,A}.
\]
From this point onward define

\[
 \omega_{X,A}:=\omega_{X/S,A},\qquad \omega_X:=\omega_{X/S,A}
\]

when \(S\) and \(A\) are fixed.  Thus every occurrence of \(\omega_X\) is
an abbreviation for the previously defined relative object.
\end{definition}

\begin{definition}[Dualizing trace]
Let \(f:X\to Y\) be separated and of finite type over \(S\).  The counit of
the typed adjunction

\[
 Rf_!:D(X,A)\rightleftarrows D(Y,A):Rf^!
\]

has component

\[
 \epsilon^!_{f,\omega_Y}:Rf_!Rf^!\omega_{Y,A}\longrightarrow\omega_{Y,A}.
\]

Define the dualizing trace to be

\[
\begin{aligned}
 \Tr_{f,A}:Rf_!\omega_{X,A}
 &\xrightarrow{Rf_!(\kappa_{f,A})}
 Rf_!Rf^!\omega_{Y,A}\\
 &\xrightarrow{\epsilon^!_{f,\omega_Y}}\omega_{Y,A}.
\end{aligned}
\]

It is a morphism in \(D(Y,A)\).  If \(S\) and \(A\) are fixed, write
\(\Tr_f\) for \(\Tr_{f,A}\).
\end{definition}

\begin{theorem}[Finiteness of the dualizing object]
If \(X\to S\) is separated and of finite type, then

\[
 \omega_{X/S,A}\in\Dbc(X,A).
\]

Its cohomological amplitude is bounded in terms of \(\dim X\), and it has the
finite-Tor property on every smooth stratum.
\end{theorem}

\begin{theorem}[Coefficient compatibility of relative dualizing objects]
For \(m\leq n\), exceptional coefficient change and the maps
\(\lambda_{n,m}\) define

\[
 \red_{n,m,X}\omega_{X/S,\Lambda_n}
 \xrightarrow{\sim}\omega_{X/S,\Lambda_m}.
\]

These maps satisfy the identity and cocycle equations and recover the
normalization of \(\omega_{X/S,\OO_E}\).  For rationalization or a finite
field extension \(A\to B\), there is a transitive isomorphism

\[
 \omega_{X/S,A}\otimes_A^LB\xrightarrow{\sim}\omega_{X/S,B}.
\]

Both comparisons commute with \(\kappa_f\), the exceptional counit, and the
typed trace \(\Tr_f\).  They therefore descend to the pair/Karoubi category
over \(\Qlb\).
\end{theorem}

\begin{definition}[Ambient Verdier dual]
At a finite stage \(A=\Lambda\), for every \(M\in D(X,\Lambda)\), define

\[
 \DD_{X/S,\Lambda}^{\mathrm{amb}}(M)
 =\RHom_{X,\Lambda}(M,\omega_{X/S,\Lambda})
 \in D(X,\Lambda).
\]

For \(u:M\to N\), define \(\DD^{\mathrm{amb}}(u)\) by
precomposition.  Thus internal Hom defines a contravariant exact
triangulated functor

\[
 \DD_{X/S,\Lambda}^{\mathrm{amb}}:
 D(X,\Lambda)^{\mathrm{op}}\longrightarrow D(X,\Lambda).
\]

At the \(\OO_E\), \(E\), and \(\Qlb\) stages, define the first Verdier dual
on \(K\in\Dctf(X,A)\) by the coefficient reconstruction

\[
 \DD_{X/S,A}K=\RHom_{X,A}(K,\omega_{X/S,A}).
\]

Whenever no ambiguity is possible, write \(\DD_X\).  The superscript
\(\mathrm{amb}\) is retained only for the all-complex finite-coefficient
functor.
\end{definition}

\begin{theorem}[Bounded constructibility of the first dual]
For every \(K\in\Dctf(X,A)\), one has

\[
 \DD_XK=\RHom_X(K,\omega_X)\in\Dbc(X,A).
\]

Consequently the ambient functor restricts to a contravariant exact functor

\[
 \DD_X:\Dctf(X,A)^{\mathrm{op}}\longrightarrow\Dbc(X,A).
\]

This statement does not yet assert that \(\DD_XK\) has finite Tor-amplitude.
\end{theorem}

\begin{proposition}[Functorial identities]
For composable maps \(K\xrightarrow uL\xrightarrow vM\),

\[
 \DD(\operatorname{id}_K)=\operatorname{id}_{\DD K},
 \qquad
 \DD(vu)=\DD(u)\DD(v).
\]

The functor sends a distinguished triangle to the oppositely oriented
distinguished triangle and satisfies
\(\DD(K[1])\cong(\DD K)[-1]\).
\end{proposition}

\begin{definition}[Ambient biduality morphism]
At a finite stage, for \(M\in D(X,\Lambda)\), define

\[
 \delta_{X,M}:M\longrightarrow
 \DD_X^{\mathrm{amb}}\DD_X^{\mathrm{amb}}M
\]

as the tensor--Hom adjoint of the evaluation morphism

\[
 M\otimes^L\RHom_X(M,\omega_X)\longrightarrow\omega_X.
\]
Both duals in this finite-level target are formed in the ambient derived
category, without a constructibility or finite-Tor prerequisite.  At an adic
or field stage, for \(K\in\Dctf(X,A)\) such that
\(\DD_XK\in\Dctf(X,A)\), define \(\delta_{X,K}\) as the unique coefficient
reconstruction of these finite-level maps.  No second dual, and hence no
biduality map, is asserted there before this finite-Tor condition is known.
\end{definition}

\begin{lemma}[Naturality of the ambient biduality morphism]
At finite coefficients, for every \(u:M\to N\) in \(D(X,\Lambda)\), the
square

\[
\begin{CD}
M @>{\delta_{X,M}}>> \DD_X^{\mathrm{amb}}\DD_X^{\mathrm{amb}}M\\
@V{u}VV @VV{\DD_X^{\mathrm{amb}}\DD_X^{\mathrm{amb}}(u)}V\\
N @>{\delta_{X,N}}>> \DD_X^{\mathrm{amb}}\DD_X^{\mathrm{amb}}N
\end{CD}
\]

commutes.  At the adic and field stages the same square is defined and
commutes when \(M,N,\DD_XM,\DD_XN\) all lie in \(\Dctf(X,A)\).
\end{lemma}

\begin{definition}[Predual-finite objects]
Define \(D_{\mathrm{pre}}(X,A)\) to be the full subcategory of
\(\Dctf(X,A)\) consisting of the objects \(K\) for which

\[
 \DD_XK=\RHom_X(K,\omega_X)\in\Dctf(X,A).
\]

This condition only asserts that the first dual belongs again to the
constructible finite-Tor subcategory.  It contains no biduality claim; the
second dual and \(\delta_{X,K}\) are then defined in \(D(X,A)\) by the
coefficient-stage rule above.
\end{definition}

\begin{definition}[Verdier-reflexive objects]
Define \(D_{\mathrm{ref}}(X,A)\) to be the full subcategory of
\(D_{\mathrm{pre}}(X,A)\) consisting of objects for which the already defined
map

\[
 \delta_{X,K}:K\longrightarrow\DD_X\DD_XK
\]

is an isomorphism.  Thus predual finiteness is checked before reflexivity.
\end{definition}

\begin{proposition}[Closure of the two finiteness domains]
The subcategories \(D_{\mathrm{pre}}(X,A)\) and
\(D_{\mathrm{ref}}(X,A)\) are thick: each contains zero and is closed under
shifts, cones, finite direct sums, and direct summands.

For an open--closed decomposition \(X=U\sqcup Z\), membership in either
subcategory is local provided the typed open and closed duality exchanges in
Theorem~\ref{thm:open-closed-exchange} are isomorphisms.
\end{proposition}

\section{Local biduality and constructible devissage}

\begin{definition}[Lisse perfect complex]
An object \(K\in D(U,A)\) is a lisse perfect complex if there is an etale
cover \((V_i\to U)\) such that every restriction \(K|_{V_i}\) is
quasi-isomorphic to the constant-sheaf realization of a bounded complex of
finite projective \(A\)-modules.  Equivalently, \(K\) is locally constant as
an object of \(D(A)\) and its coefficient complex is perfect.  Its
cohomology sheaves are lisse and finite, but their individual coefficient
modules are not required to be projective or perfect.
\end{definition}

\begin{theorem}[Smooth-stratum formula]
Let \(i:U\to S\) be smooth of pure relative dimension \(d\).  Smooth purity
gives

\[
 \omega_{U/S}\xrightarrow{\sim}Li^*\omega_S(d)[2d].
\]

If \(K\) is lisse and perfect on \(U\), then

\[
 \DD_UK\cong K^\vee\otimes^L\omega_{U/S},
\]

and the dual-basis coevaluation defines an inverse
\(\DD_U\DD_UK\to K\) to \(\delta_{U,K}\).
\end{theorem}

\begin{proposition}[Basis independence]
The dual-basis inverse on a lisse perfect complex is independent of every
local trivialization.  It is natural in the complex, compatible with shifts
and cones, and compatible with etale pullback.
\end{proposition}

\begin{definition}[Duality stratification over the base]
A duality stratification of \(K\in\Dctf(X,A)\) is a finite decomposition

\[
 X=\coprod_{r=1}^mX_r
\]

by locally closed subschemes smooth over \(S\), such that the restricted
complex \(K|_{X_r}\) is lisse perfect on every \(X_r\).  The strata may be
ordered so that
\(U_r=\coprod_{i\leq r}X_i\) is open in \(X\).
\end{definition}

\begin{theorem}[Existence of duality stratifications over a perfect field]
Let \(S=\operatorname{Spec}k\), where \(k\) is perfect, and let \(X\) be of
finite type over \(k\).  Every object of \(\Dctf(X,A)\) admits a duality
stratification over \(k\).  Any two such stratifications admit a common
finite refinement by smooth locally closed \(k\)-subschemes.
\end{theorem}

\begin{theorem}[Open--closed localization]
Let \(j:U\hookrightarrow X\) be a quasi-compact open immersion and let
\(i:Z\hookrightarrow X\) be its closed complement.  The functors occurring
below have types

\[
 j_!,Rj_*:D(U,A)\longrightarrow D(X,A),\qquad
 j^*:D(X,A)\longrightarrow D(U,A),
\]

\[
 i_*=Ri_*:D(Z,A)\longrightarrow D(X,A),\qquad
 i^*,Ri^!:D(X,A)\longrightarrow D(Z,A).
\]

For \(K\in\Dbc(X,A)\), there are distinguished triangles in \(D(X,A)\)

\[
 j_!j^*K\longrightarrow K\longrightarrow i_*i^*K
 \longrightarrow j_!j^*K[1],
\]

\[
 i_*Ri^!K\longrightarrow K\longrightarrow Rj_*j^*K
 \longrightarrow i_*Ri^!K[1].
\]

All six displayed functors preserve bounded constructibility on the objects
in these triangles; every finite-Tor restriction is used only under its
separate preservation theorem.
\end{theorem}

\begin{theorem}[Duality across an open--closed decomposition]
\label{thm:open-closed-exchange}
Let \(j:U\hookrightarrow X\) and \(i:Z\hookrightarrow X\) be as in the
preceding theorem.  For \(K\in\Dctf(U,A)\) and \(L\in\Dctf(Z,A)\), whenever
the displayed source objects lie in the finite-Tor domain, there are natural
isomorphisms in \(D(X,A)\)

\[
 \DD_X(j_!K)\xrightarrow{\sim}Rj_*(\DD_UK),
 \qquad
 \DD_X(Rj_*K)\xrightarrow{\sim}j_!(\DD_UK),
\]

\[
 \DD_X(i_*L)\xrightarrow{\sim}i_*(\DD_ZL).
\]

For \(M\in\Dctf(X,A)\) with \(Ri^!M\in\Dctf(Z,A)\), there is a natural
isomorphism in \(D(Z,A)\)

\[
 i^*(\DD_XM)\xrightarrow{\sim}\DD_Z(Ri^!M).
\]

Under these isomorphisms, applying \(\DD_X\) exchanges the two localization
triangles whenever their three terms belong to \(D_{\mathrm{pre}}\).  If
they belong to \(D_{\mathrm{ref}}\), the three biduality isomorphisms form a
morphism from the original triangle to its double dual.
\end{theorem}

\begin{proposition}[One-stratum biduality]
If \(K\) is supported on one smooth stratum and is lisse perfect there, then
\(\delta_{X,K}\) is an isomorphism.
\end{proposition}

\begin{proposition}[Induction step for biduality]
In an open--closed localization triangle, if the biduality morphisms for the
open restriction and the closed restriction are isomorphisms, then the
biduality morphism for the middle object is an isomorphism.
\end{proposition}

\begin{theorem}[Stratified Verdier biduality]
Let \(K\in\Dctf(X,A)\) be supplied with a duality stratification over \(S\).
Then \(\DD_XK\in\Dctf(X,A)\), and the natural morphism

\[
 \delta_{X,K}:K\xrightarrow{\sim}\DD_X\DD_XK
\]

is an isomorphism.  Equivalently, \(K\in D_{\mathrm{ref}}(X,A)\).
\end{theorem}

\begin{corollary}[Verdier biduality over a perfect field]
Let \(S=\operatorname{Spec}k\) with \(k\) perfect, and let \(X\) be of finite
type over \(k\).  Then every object of \(\Dctf(X,A)\) is Verdier-reflexive,
and \(\DD_X\) is an anti-equivalence

\[
 \DD_X:\Dctf(X,A)^{\mathrm{op}}\xrightarrow{\sim}\Dctf(X,A)
\]

whose quasi-inverse is itself.
Equivalently,
\[
 D_{\mathrm{ref}}(X,A)=D_{\mathrm{pre}}(X,A)=\Dctf(X,A).
\]
\end{corollary}

\section{Relative duality and exchange transformations}

\begin{convention}[Biduality scope]
In this section, an invocation of biduality assumes either that the object in
question is equipped with a duality stratification, or that the schemes are
of finite type over a perfect field.  Relative-duality statements not using
biduality require the ambient exceptional adjunction together with the
separate preservation and exchange prerequisites stated in Section~1.
\end{convention}

\begin{definition}[Mixed ordinary/compact tensor morphism]
Let \(f:X\to Y\) be separated and of finite type, let
\(K\in\Dctf(X,A)\), and let \(M\in\Dbc(X,A)\).  Assume the compact-support
projection comparison is invertible on the displayed objects.  Define

\[
 \mu_{f,K,M}:Rf_!K\otimes^L_{A_Y}Rf_*M
 \longrightarrow Rf_!(K\otimes^L_{A_X}M)
\]

as the typed composite

\[
\begin{aligned}
 Rf_!K\otimes^LRf_*M
 &\xrightarrow{\operatorname{PF}_{f,K,Rf_*M}^{-1}}
 Rf_!\bigl(K\otimes^LLf^*Rf_*M\bigr)\\
 &\xrightarrow{Rf_!(\operatorname{id}_K\otimes\epsilon^*_{f,M})}
 Rf_!(K\otimes^LM),
\end{aligned}
\]

where
\(\epsilon^*_{f,M}:Lf^*Rf_*M\to M\) is the counit of
\(Lf^*\dashv Rf_*\).  The map \(\mu_f\) is not the ordinary cup product and
is part of the data needed to type relative duality.
\end{definition}

\begin{proposition}[Coherence of the mixed tensor morphism]
The maps \(\mu_{f,K,M}\) are natural in \(K\) and \(M\), unital for the
identity morphism, and associative for composition after inserting the
composition isomorphisms for \(Rf_!\) and \(Rf_*\).  They commute with every
base-change and coefficient-change square for which the corresponding
projection, ordinary-counit, and direct-image comparisons are invertible.
\end{proposition}

\begin{definition}[Relative duality morphism]
\label{def:relative-duality}
Let \(f:X\to Y\) be separated and of finite type.  For
\(K\in\Dctf(X,A)\) and \(L\in\Dbc(Y,A)\), the projection formula and the
exceptional counit define

\[
 \operatorname{RD}_{f,K,L}:Rf_*\RHom_X(K,Rf^!L)
 \longrightarrow\RHom_Y(Rf_!K,L).
\]

It is the morphism adjoint to

\[
\begin{aligned}
 Rf_!K\otimes^LRf_*\RHom_X(K,Rf^!L)
 &\xrightarrow{\mu_{f,K,\RHom_X(K,Rf^!L)}}
 Rf_!\bigl(K\otimes^L\RHom_X(K,Rf^!L)\bigr)\\
 &\xrightarrow{Rf_!(\operatorname{ev}_{K,Rf^!L})}
 Rf_!Rf^!L
 \xrightarrow{\epsilon^!_{f,L}}L,
\end{aligned}
\]

where \(\epsilon^!_{f,L}\) is the counit of \(Rf_!\dashv Rf^!\).
\end{definition}

\begin{theorem}[Relative Verdier duality]
For \(K\in\Dctf(X,A)\) and \(L\in\Dbc(Y,A)\), the relative duality morphism
is an isomorphism

\[
 Rf_*\RHom_X(K,Rf^!L)
 \xrightarrow{\sim}\RHom_Y(Rf_!K,L).
\]

It is natural in both variables and compatible with composition of
separated finite-type morphisms.
\end{theorem}

\begin{corollary}[Verdier exchange for compact support]
Using \(Rf^!\omega_Y\cong\omega_X\), relative duality gives

\[
 Rf_*\DD_XK\xrightarrow{\sim}\DD_YRf_!K.
\]

After biduality, there is also a canonical isomorphism

\[
 Rf_!\DD_XK\xrightarrow{\sim}\DD_YRf_*K.
\]
\end{corollary}

\begin{corollary}[Proper exchange]
If \(f\) is proper, then \(Rf_!=Rf_*\), and both preceding comparisons are
the same natural isomorphism

\[
 Rf_*\DD_XK\xrightarrow{\sim}\DD_YRf_*K.
\]
\end{corollary}

\begin{corollary}[Open exchange]
For an open immersion \(j:U\hookrightarrow X\), relative duality specializes
to

\[
 Rj_*\DD_UK\xrightarrow{\sim}\DD_Xj_!K,
 \qquad
 j_!\DD_UK\xrightarrow{\sim}\DD_XRj_*K.
\]
\end{corollary}

\begin{corollary}[Smooth exchange]
If \(f:X\to Y\) is smooth of pure relative dimension \(d\), then smooth
purity identifies \(Rf^!\omega_Y\) with
\(Lf^*\omega_Y(d)[2d]\).  For \(K\in\Dctf(Y,A)\), there is a canonical
isomorphism

\[
 \DD_X(Lf^*K)
 \xrightarrow{\sim}Lf^*(\DD_YK)(d)[2d].
\]
\end{corollary}

\section{Projection, evaluation, and pairings}

\begin{definition}[Verdier evaluation pairing]
For \(K\in\Dctf(X,A)\), define

\[
 K\otimes^L\DD_XK\xrightarrow{\operatorname{ev}_K}\omega_X.
\]

Applying \(Rf_!\) and the projection formula gives, for
\(f:X\to Y\), a pairing

\[
\begin{aligned}
 Rf_!K\otimes^LRf_*\DD_XK
 &\xrightarrow{\mu_{f,K,\DD_XK}}
 Rf_!(K\otimes^L\DD_XK)\\
 &\xrightarrow{Rf_!(\operatorname{ev}_K)}Rf_!\omega_X
 \xrightarrow{\Tr_f} \omega_Y.
\end{aligned}
\]
\end{definition}

\begin{theorem}[Perfectness of the relative pairing]
The tensor--Hom adjoint of the preceding pairing is the Verdier exchange
isomorphism

\[
 Rf_*\DD_XK\xrightarrow{\sim}\DD_YRf_!K.
\]

Its adjoint in the other variable is the bidual form

\[
 Rf_!K\xrightarrow{\sim}\DD_YRf_*\DD_XK.
\]
\end{theorem}

\begin{definition}[A duality base-change square]
Let \(f:X\to Y\) be separated and of finite type, let \(u\) and \(v\) be
quasi-compact and quasi-separated, and let

\[
\begin{CD}
X' @>{u}>> X\\
@V{f'}VV @VV{f}V\\
Y' @>{v}>> Y,
\end{CD}
\]

be Cartesian.  For \(K\in\Dctf(X,A)\), the compact-support exchange has
type

\[
 \operatorname{BC}_{!,K}:Lv^*Rf_!K
 \longrightarrow Rf'_!Lu^*K
\]

in \(D(Y',A)\).  If \(f\) is proper and \(M\in\Dbc(X,A)\), proper base
change has type

\[
 \operatorname{BC}_{*,M}:Lv^*Rf_*M
 \longrightarrow Rf'_*Lu^*M
\]

in \(D(Y',A)\).  For \(N\in D(Y',A)\), the right-adjoint mate of
\(\operatorname{BC}_!\) has type

\[
 Ru_*Rf'^!N\longrightarrow Rf^!Rv_*N
\]

in \(D(X,A)\).  If \(\operatorname{BC}_!\) is invertible on the relevant
domain, denote the inverse mate by

\[
 \operatorname{BC}^{!,*}_N:Rf^!Rv_*N
 \xrightarrow{\sim}Ru_*Rf'^!N.
\]

For \(L\in\Dctf(Y,A)\), an exceptional pullback exchange is a specified map

\[
 \operatorname{Ex}^!_{f,v,L}:Lu^*Rf^!L
 \longrightarrow Rf'^!Lv^*L
\]

in \(D(X',A)\).  Call the square a \emph{duality base-change square for
\((K,L)\)} if \(\operatorname{BC}_{!,K}\) and
\(\operatorname{Ex}^!_{f,v,L}\) are isomorphisms, all four pullbacks have the
required finite-Tor landing properties, and the canonical maps

\[
 Lu^*\RHom_X(K,Rf^!L)\longrightarrow
 \RHom_{X'}(Lu^*K,Lu^*Rf^!L)
\]

and

\[
 Lv^*\RHom_Y(Rf_!K,L)\longrightarrow
 \RHom_{Y'}(Lv^*Rf_!K,Lv^*L)
\]

are isomorphisms.  If the relative-duality rectangle below is invoked, the
definition additionally requires \(f\) proper and
\(\operatorname{BC}_{*,M}\) invertible for

\[
 M=\RHom_X(K,Rf^!L)\in\Dbc(X,A).
\]
\end{definition}

\begin{theorem}[Typed base-change rectangle for relative duality]
Let the preceding square be a duality base-change square for
\(K\in\Dctf(X,A)\) and \(L\in\Dctf(Y,A)\), and assume \(f\) proper.  The
following square commutes:

\[
\begin{CD}
Lv^*Rf_*\RHom_X(K,Rf^!L)
 @>{Lv^*(\operatorname{RD}_{f,K,L})}>>
Lv^*\RHom_Y(Rf_!K,L)\\
@V{\sim}VV @VV{\sim}V\\
Rf'_*\RHom_{X'}(Lu^*K,Rf'^!Lv^*L)
 @>{\operatorname{RD}_{f',Lu^*K,Lv^*L}}>>
\RHom_{Y'}(Rf'_!Lu^*K,Lv^*L).
\end{CD}
\]

The left vertical arrow is the composite of
\(\operatorname{BC}_{*,\RHom_X(K,Rf^!L)}\), the first internal-Hom pullback
comparison, and \(\operatorname{Ex}^!_{f,v,L}\).  The right vertical arrow is
the second internal-Hom pullback comparison followed contravariantly by
\(\operatorname{BC}_{!,K}^{-1}\).  Thus both vertical arrows are morphisms in
\(D(Y',A)\).  The rectangle is transitive under pasting of duality
base-change squares.
\end{theorem}

\begin{theorem}[Typed base-change rectangle for the evaluation pairing]
Let \(K\in D_{\mathrm{pre}}(X,A)\), assume that
\(Lu^*K\in D_{\mathrm{pre}}(X',A)\), and assume in addition that

\[
 Lu^*\omega_X\xrightarrow{\sim}\omega_{X'},\qquad
 Lv^*\omega_Y\xrightarrow{\sim}\omega_{Y'}
\]

are specified isomorphisms in \(D(X',A)\) and \(D(Y',A)\), respectively,
induced by the chosen compatible base dualizing objects.  Require them to
commute with \(\kappa_f\) and the counits defining \(\Tr_f\).  Under
\(\operatorname{BC}_!\), \(\operatorname{BC}_*\), and the
induced comparison \(Lu^*\DD_XK\xrightarrow{\sim}\DD_{X'}Lu^*K\), pullback
of the typed pairing

\[
 Rf_!K\otimes^LRf_*\DD_XK\longrightarrow\omega_Y
\]

is the pairing

\[
 Rf'_!Lu^*K\otimes^LRf'_*\DD_{X'}Lu^*K\longrightarrow\omega_{Y'}.
\]

These pairing rectangles are transitive under pasting.
\end{theorem}

\section{Coefficient rectangles}

\begin{theorem}[Finite reduction of internal Hom]
Let \(m\leq n\).  For \(K\in\Dctf(X,\Lambda_n)\) and
\(L\in\Dbc(X,\Lambda_n)\), the canonical map is an isomorphism

\[
 \red_{n,m,X}\RHom_{X,\Lambda_n}(K,L)
 \xrightarrow{\sim}
 \RHom_{X,\Lambda_m}(\red K,\red L).
\]

For \(M\in\Dbc(X,\Lambda_n)\), evaluation is compatible with the comparison:

\[
\begin{CD}
\red K\otimes^L_{\Lambda_{m,X}}\red\RHom_{\Lambda_n}(K,M)
 @>{\operatorname{id}\otimes\sim}>>
\red K\otimes^L\RHom_{\Lambda_m}(\red K,\red M)\\
@V{\red(\operatorname{ev}_{K,M})}VV
 @VV{\operatorname{ev}_{\red K,\red M}}V\\
\red M @= \red M.
\end{CD}
\]

The corresponding currying rectangle commutes for
\(P,K\in\Dctf(X,\Lambda_n)\) and \(M\in\Dbc(X,\Lambda_n)\).  Compatibility
with a six-operation functor is used only after invoking that functor's own
coefficient-change theorem.
\end{theorem}

\begin{theorem}[Reduction of dualizing objects and Verdier duals]
Let \(m\leq n\) and let \(K\in\Dctf(X,\Lambda_n)\).  
There are canonical isomorphisms

\[
 \red_{n,m,X}\omega_{X,\Lambda_n}
 \xrightarrow{\sim}\omega_{X,\Lambda_m},
\]

\[
 \red_{n,m,X}\DD_{X,\Lambda_n}K
 \xrightarrow{\sim}\DD_{X,\Lambda_m}(\red K).
\]

For \(K,L\in\Dctf(X,\Lambda_n)\) and every \(u:K\to L\), the rectangle

\[
\begin{CD}
\red\DD_{\Lambda_n}L @>{\sim}>> \DD_{\Lambda_m}(\red L)\\
@V{\red(\DD u)}VV @VV{\DD(\red u)}V\\
\red\DD_{\Lambda_n}K @>{\sim}>> \DD_{\Lambda_m}(\red K)
\end{CD}
\]

commutes.  The analogous rectangles commute for \(\delta\), relative
duality, proper exchange, open exchange, smooth exchange, and evaluation
pairings whenever every object involved lies in the displayed finite-Tor
domains.
\end{theorem}

\begin{theorem}[Adic reconstruction of Verdier duality]
For \(K=(K_n)\in\Dctf(X,\OO_E)\), the system
\((\DD_{X,\Lambda_n}K_n)_n\) is normalized and defines the first dual
\(\DD_{X,\OO_E}K\in\Dbc(X,\OO_E)\).  It lies in
\(\Dctf(X,\OO_E)\) precisely under the common finite-Tor bound required in
the definition of \(D_{\mathrm{pre}}(X,\OO_E)\).  Only under that condition
do the finite-level maps \(\delta_{X,K_n}\) define the adic biduality map.

If the \(K_n\) have compatible duality stratifications with common bounds,
or if \(X\) is of finite type over the fixed perfect field, every finite-level
biduality map is an isomorphism and the reconstructed adic map is an
isomorphism.  Relative-duality, exchange, and pairing isomorphisms reconstruct
only when all their displayed objects satisfy the same uniform landing
conditions.  Evaluation at level \(n\) returns each finite-level map exactly.
\end{theorem}

\begin{theorem}[Scalar extension]
Let \(A\to B\) be rationalization or a finite extension of coefficient
fields.  For \(K\in\Dctf(X,A)\), the comparison map is an isomorphism

\[
 \DD_{X,A}K\otimes_A^LB
 \xrightarrow{\sim}
 \DD_{X,B}(K\otimes_A^LB).
\]

It is transitive in towers and independent of an integral lattice.  If
\(K\in D_{\mathrm{pre}}(X,A)\), it transports the defined biduality map; if
\(K\in D_{\mathrm{ref}}(X,A)\), it transports the biduality isomorphism.
Exchange, relative-duality, and pairing diagrams are transported when all
their separately stated finite-Tor and base-change hypotheses hold after
extension to \(B\).
\end{theorem}

\begin{theorem}[\(\Qlb\)-coefficient duality]
Represent \(K\in\Dctf(X,\Qlb)\) by a triple \((E,K_E,e)\).  The first dual
is represented by \((E,\DD_{X,E}K_E,\DD_{X,E}(e))\), where contravariance
reverses multiplication of endomorphisms and preserves idempotence.  It is
independent of every further finite extension.  A second dual and its
biduality map are defined only when the represented first dual belongs to
\(\Dctf\); the biduality map is an isomorphism only when the representing
object lies in \(D_{\mathrm{ref}}(X,E)\).  Relative duality, exchanges, and
pairings descend to \(D(X,\Qlb)=\operatorname{Kar}(\mathcal C_X)\) under
their separately stated landing and coefficient-change hypotheses.
\end{theorem}

\end{document}
